\documentclass[11pt,reqno]{amsart}
\usepackage{amsmath,amssymb,amsthm,geometry,mathtools}
\geometry{a4paper,margin=1.1in}

\newtheorem{theorem}{Theorem}[section]
\newtheorem{lemma}[theorem]{Lemma}
\newtheorem{proposition}[theorem]{Proposition}
\newtheorem{remark}[theorem]{Remark}

\newcommand{\e}{\mathrm{e}}
\newcommand{\m}{\mathfrak{m}}
\newcommand{\M}{\mathfrak{M}}
\newcommand{\R}{\mathbb{R}}
\newcommand{\Z}{\mathbb{Z}}
\newcommand{\N}{\mathbb{N}}

\title{An Analytic Evaluation of Estimate (K) and the Circle Method Defect for OEIS A306477}
\author{Raj123-0}
\date{\today}

\begin{document}

\maketitle

\begin{abstract}
We evaluate the minor arc $L^2$ bound (Estimate K) for the degree-8 Weyl sum $f_8(\alpha) = \sum_{1 \le z \le N^{1/8}} \e\left(\alpha \binom{z+7}{8}\right)$. We prove that the total $L^2$ energy satisfies $\|f_8\|_{L^2([0,1])}^2 = \lfloor N^{1/8} \rfloor \le N^{1/4 - 1/8}$, establishing Estimate (K) with exact parameter $\delta = 1/8$. However, we demonstrate that applying Estimate (K) via Cauchy--Schwarz produces an upper bound of $N^{23/48}$, which exceeds the singular integral main term $J(n) \asymp n^{1/24} = n^{2/48}$ by an exact dimensional deficit of $\Delta = 7/16 = 0.4375$. We formulate the required joint bilinear cancellation condition.
\end{abstract}

\section{Exact Evaluation of Estimate (K)}

Let $P = \lfloor N^{1/8} \rfloor$, and let $f_8(\alpha) = \sum_{z=1}^P \e\left(\alpha \binom{z+7}{8}\right)$.

\begin{theorem}[Orthogonality and $L^2$ Bound]
For any measurable subset $\m \subseteq [0, 1]$ and any $N \ge 1$:
\begin{equation}
    \int_{\m} |f_8(\alpha)|^2 \, d\alpha \le \int_0^1 |f_8(\alpha)|^2 \, d\alpha = P \le N^{1/8}.
\end{equation}
Setting $\delta = 1/8 = 0.125$ and $C = 1$, we have:
\begin{equation}
    \int_{\m} |f_8(\alpha)|^2 \, d\alpha \le 1 \cdot N^{1/4 - \delta}.
\end{equation}
\end{theorem}

\begin{proof}
Expanding the square and integrating term-by-term:
\begin{equation}
    \int_0^1 |f_8(\alpha)|^2 \, d\alpha = \sum_{1 \le z_1, z_2 \le P} \int_0^1 \e\left(\alpha \left(\binom{z_1+7}{8} - \binom{z_2+7}{8}\right)\right) \, d\alpha.
\end{equation}
By the orthogonality relation $\int_0^1 \e(m\alpha) \, d\alpha = \delta_{m,0}$, the integral is non-zero if and only if $\binom{z_1+7}{8} = \binom{z_2+7}{8}$. Since $z \mapsto \binom{z+7}{8}$ is strictly increasing for $z \ge 1$, this forces $z_1 = z_2$. There are exactly $P$ such diagonal solutions and zero off-diagonal solutions.
\end{proof}

\section{The Circle Method Dimensional Obstruction}

When inserting this bound into the 4-variable representation integral:
\begin{equation}
    \left| \int_{\m} f_2 f_4 f_6 f_8 \, \e(-n\alpha) \, d\alpha \right| \le \|f_8\|_{L^2(\m)} \|f_2 f_4 f_6\|_{L^2([0,1])} \le N^{1/16} \cdot N^{5/12} = N^{23/48}.
\end{equation}
Meanwhile, the major arc main term is:
\begin{equation}
    \int_{\M} f_2 f_4 f_6 f_8 \, \e(-n\alpha) \, d\alpha = \mathfrak{S}(n) J(n) + o(n^{1/24}) \asymp n^{1/24} = n^{2/48}.
\end{equation}
The uncoupled minor arc bound exceeds the main term by:
\begin{equation}
    \Delta = \frac{23}{48} - \frac{2}{48} = \frac{21}{48} = \frac{7}{16} = 0.4375.
\end{equation}
Hence, uncoupling $f_8$ from $(f_2, f_4, f_6)$ cannot establish $r(n) > 0$. A proof requires proving joint phase cancellation directly on the 4-fold product.

\end{document}